\chapter{Configuración experimental}
\label{app:configuracion}

Este apéndice recoge los parámetros principales de la configuración experimental. Se incluyen los valores necesarios para interpretar los resultados y comparar los enfoques CNN y VLM/ICL.

\section{Conjunto sintético QRS/ST}

\begin{table}[H]
  \centering
  \small
  \begin{tabularx}{0.94\textwidth}{p{4.4cm}X}
    \toprule
    Propiedad & Valor \\
    \midrule
    Fuente & Simulador QRS/ST descrito en el capítulo~\ref{chap:simulador} \\
    Pacientes & 100 \\
    Imágenes únicas & 100 \\
    Predicciones agregadas & 300, correspondientes a tres semillas \\
    Familias fuente & Cinco familias morfológicas descritas en el capítulo~\ref{chap:simulador} \\
    Pacientes por familia & 20 \\
    Derivaciones renderizadas & V1 \\
    Frecuencia de muestreo & 500 Hz \\
    Duración del latido & 900 ms \\
    Resolución de guardado & 130 ppp \\
    Positivos derivados & 20 \\
    Negativos o incompletos & 80 \\
    Valores de \(K\) & \(K=\{2,4,8,16,32\}\) \\
    Semillas & 42, 123 y 2026 \\
    \bottomrule
  \end{tabularx}
  \caption{Configuración del conjunto sintético final.}
  \label{tab:config_sintetico_final}
\end{table}

La etiqueta positiva se obtiene cuando RBBB, elevación del ST e inversión de T están presentes de forma simultánea. La decisión por paciente utiliza únicamente la imagen V1 y aplica después el umbral seleccionado.

\section{Cohorte del HUCA}

\begin{table}[H]
  \centering
  \small
  \begin{tabularx}{0.94\textwidth}{p{4.4cm}X}
    \toprule
    Propiedad & Valor \\
    \midrule
    Fuente & Hospital Universitario Central de Asturias (HUCA), recurso Brugada HUCA 1.0.0 \\
    Pacientes brutos & 363 \\
    Pacientes excluidos & 46 por patrón basal limítrofe \\
    Pacientes incluidos & 317 \\
    Imágenes únicas incluidas & 317 \\
    Predicciones agregadas & 951, correspondientes a tres semillas \\
Referencias positivas & 116 \\
Referencias negativas & 201 \\
    Derivaciones & V1 \\
    Ventana & 300 ms antes de R y 600 ms después de R \\
    Frecuencia de origen & 100 Hz \\
    Resolución de guardado & 130 ppp \\
    Separación & Leave-one-out por paciente \\
    \bottomrule
  \end{tabularx}
  \caption{Configuración de la cohorte procesada del HUCA.}
  \label{tab:config_huca_final}
\end{table}

La cohorte procesada del HUCA no incluye anotaciones independientes de los tres hallazgos morfológicos simulados. La evaluación emplea una referencia clínica por paciente. Esta variable, almacenada en \texttt{clinical\_brugada}, procede de una revisión especializada.

\section{Configuración CNN}

\begin{table}[H]
  \centering
  \small
  \begin{tabularx}{0.94\textwidth}{p{4.4cm}X}
    \toprule
    Propiedad & Valor \\
    \midrule
    Arquitectura & ResNet18 \\
    Pesos iniciales & Inicialización estándar de ResNet18 \\
    Entrada & RGB, 224 por 224 \\
    Salida & Tres hallazgos morfológicos \\
    Pérdida & Entropía cruzada binaria multietiqueta \\
    Optimizador & AdamW \\
    Tasa de aprendizaje & \(3\times10^{-4}\) \\
    Lote & 32 \\
    Épocas máximas & 20 \\
    Estrategia de umbral & Selección en validación mediante Balanced Accuracy \\
    Métrica principal & Balanced Accuracy por paciente \\
    Grad-CAM & Límite de 12 paneles por ejecución sintética y 6 por ejecución sobre el HUCA \\
    \bottomrule
  \end{tabularx}
  \caption{Hiperparámetros de la CNN final.}
  \label{tab:config_cnn_appendix}
\end{table}

\section{Configuración de adaptación de dominio}

\begin{table}[H]
  \centering
  \small
  \begin{tabularx}{0.94\textwidth}{p{4.4cm}X}
    \toprule
    Propiedad & Valor \\
    \midrule
    Origen etiquetado & Conjunto sintético QRS/ST \\
    Destino no etiquetado & Cohorte del HUCA \\
    Métodos & SSL, CORAL, MMD y DANN \cite{chen2020simclr,sun2016deepcoral,long2015dan,ganin2016} \\
    Valores de \(K\) & \(K=16\) y \(K=32\) \\
    Semillas & 42, 123 y 2026 \\
    Peso de adaptación & 0,1 \\
    Escalas MMD & 0,5, 1,0, 2,0 y 4,0 \\
    SSL & SimCLR durante 3 épocas \\
    Política de destino & Transductiva compartida, sin etiquetas destino en la pérdida supervisada \\
    \bottomrule
  \end{tabularx}
  \caption{Configuración de los experimentos de adaptación de dominio.}
  \label{tab:config_domain_appendix}
\end{table}

\clearpage
\section{Configuración VLM/ICL}

El enfoque VLM/ICL usa los mismos pacientes reservados, valores de \(K\) y métricas por paciente que la CNN. La evaluación final corresponde a Gemma~4, servido mediante vLLM en infraestructura universitaria, con temperatura 0 y salida estructurada.

\begin{table}[H]
  \centering
  \small
  \begin{tabularx}{0.94\textwidth}{p{4.4cm}X}
    \toprule
    Elemento & Valor \\
    \midrule
    Modelo informado & \texttt{google/gemma-4-E4B-it} \\
    Servicio & vLLM en infraestructura universitaria \\
    Entrada & Imagen ECG y demostraciones imagen-etiqueta \\
    Valores de \(K\) & \(K=\{0,2,4,8,16,32\}\) \\
    Condiciones informadas & Zero-shot, ICL estándar e ICL balanceado \\
    Contextos sobre la cohorte del HUCA & Demostraciones de la propia cohorte o del simulador \\
    Salida morfológica & RBBB, elevación del ST e inversión de T \\
    Salida clínica & Caso Brugada confirmado frente a normal \\
    Métrica principal & Balanced Accuracy por paciente \\
    \bottomrule
  \end{tabularx}
  \caption{Configuración resumida del enfoque VLM/ICL.}
  \label{tab:config_vlm_appendix}
\end{table}

La sección~\ref{chap:vlm} describe este enfoque y el capítulo de resultados desarrolla su comparación con la CNN.
